\documentclass[11pt]{article}
\usepackage[margin=1in]{geometry}
\usepackage{graphicx}
\usepackage{booktabs}
\usepackage{float}

\title{Advanced Econometric Methods:sample}
\author{Student ID: Your Student ID\\Your Name}
\date{}

% Keep the template compilable before Stata output exists.
\newcommand{\resultfile}[2]{%
  \IfFileExists{#1}{\input{#1}}{\fbox{\parbox{0.9\linewidth}{#2}}}%
}

\begin{document}
\maketitle

% Rename every section to match the current assignment.
% Add or remove sections when the number of questions changes.

\section*{[Question number and short title]}
Answer the question directly. Then present the requested result and explain its direction, magnitude, and substantive meaning.

% Example of a code-generated table. Change or remove as needed.
\begin{table}[H]
  \centering
  \caption{[Descriptive table title]}
  \label{tab:main}
  \resultfile{../output/tables/table_main.tex}{Run code/analysis.do to create output/tables/table\_main.tex.}
\end{table}
The main result is ...

\section*{[Next question number and short title]}
Answer and interpretation: ...

% Example of a code-generated figure. Change or remove as needed.
\begin{figure}[H]
  \centering
  \IfFileExists{../output/figures/figure_main.pdf}
    {\includegraphics[width=0.75\textwidth]{../output/figures/figure_main.pdf}}
    {\fbox{Run code/analysis.do to create output/figures/figure\_main.pdf.}}
  \caption{[Descriptive figure title; include units or a note if needed]}
  \label{fig:main}
\end{figure}
Figure~\ref{fig:main} shows ...

% Continue by copying the section pattern as many times as needed.

\end{document}
